\documentclass{article}
\usepackage{../../third_party/iclr2027/iclr2027_conference,times}
\input{../../third_party/iclr2027/math_commands.tex}
% Generated by scripts/build_paper_claim_summary.py; do not edit by hand.
\newcommand{\PTLPrimaryDeltaCross}{0.324175}
\newcommand{\PTLPrimaryDeltaAdj}{0.093331}
\newcommand{\PTLPrimaryDeltaAdjCI}{[0.056748, 0.134389]}
\newcommand{\PTLPrimarySystemaCross}{0.110196}
\newcommand{\PTLPrimarySystemaAdj}{0.035181}
\newcommand{\PTLPrimarySystemaAdjCI}{[0.017110, 0.057823]}
\newcommand{\PTLPrimaryRankCross}{0.328296}
\newcommand{\PTLPrimaryRankAdj}{0.073300}
\newcommand{\PTLPrimaryRankAdjCI}{[0.035638, 0.111859]}
\newcommand{\PTLDecisionRho}{0.775026}
\newcommand{\PTLDecisionRhoDisplay}{0.78}
\newcommand{\PTLDecisionRhoCI}{[0.73, 0.89]}
\newcommand{\PTLWithinMetricRho}{0.26}
\newcommand{\PTLWithinMetricRhoCI}{[0.02, 0.45]}
\newcommand{\PTLMeanFidelityShiftRhoD}{0.69}
\newcommand{\PTLMeanFidelityShiftRhoRegret}{0.77}
\newcommand{\PTLSameContextRange}{0.0012--0.0054}
\newcommand{\PTLNadigDelta}{0.176}
\newcommand{\PTLNadigSystema}{0.138}
\newcommand{\PTLNadigRank}{0.129}
\newcommand{\PTLFrangiehLabels}{243}

% Generated by scripts/build_scientific_upgrade_tables.py; do not edit by hand.
\newcommand{\PTLDecompositionIdentityError}{8.33e-17}
\newcommand{\PTLDecompositionVsCanonicalMaxDiff}{0.022}
\newcommand{\PTLComparatorKendallRho}{0.87}
\newcommand{\PTLComparatorJaccardRho}{0.84}
\newcommand{\PTLIncrementalMeanDeltaMAE}{0.013}
\newcommand{\PTLIncrementalMeanDeltaMAECI}{[-0.001, 0.027]}
\newcommand{\PTLIncrementalDAdjMetricOnlyDeltaMAE}{-0.005}
\newcommand{\PTLIncrementalDAdjMetricOnlyDeltaMAECI}{[-0.019, 0.008]}
\newcommand{\PTLIncrementalDAdjBeyondMeanDeltaMAE}{-0.005}
\newcommand{\PTLIncrementalDAdjBeyondMeanDeltaMAECI}{[-0.014, 0.003]}
\newcommand{\PTLTargetAuditRows}{360}
\newcommand{\PTLTargetAuditFalseReuse}{0.019}

\usepackage{graphicx}
\usepackage{booktabs}
\usepackage{array}
\usepackage{float}
\usepackage{placeins}
\usepackage{hyperref}
\usepackage{url}
\hypersetup{hidelinks}

\renewcommand{\tablename}{Supplementary Table}
\renewcommand{\thetable}{S\arabic{table}}
\renewcommand{\figurename}{Supplementary Figure}
\renewcommand{\thefigure}{S\arabic{figure}}
\renewcommand{\thesection}{\Alph{section}}

\title{Supplement to ``Measuring the Transport of Perturbation Reliability Rankings Across Biological Contexts''}
\author{Anonymous Authors}

\begin{document}
\maketitle

\noindent\textbf{Organization.} This appendix supplies the data and source-frozen protocol (A), estimators and uncertainty decomposition (B), extended transport and independent replication evidence (C), measurement and decision robustness (D), the prospective source-only boundary (E), additional robustness and diagnostic analyses (F), the external predictor context under matched information constraints (G), and the scientific regeneration path (H). It complements rather than repeats the main-text sequence: observed reordering $\rightarrow$ measurement-adjusted transport $\rightarrow$ independent replication $\rightarrow$ measurement boundary $\rightarrow$ decision consequence.

\section{Data and source-frozen evaluation protocol}
\label{app:data}

The primary surface is the Frangieh Perturb-CITE-seq melanoma data set \citep{frangieh2021multimodal}. We normalize each cell to a total of 10,000, apply $\log(1+x)$, and retain cells with at least 700 detected genes, 1,500 counts, and at most 25\% mitochondrial counts. After cell QC, the analysis uses the fixed 8,229-gene evaluation panel. Control, co-culture, and IFN$\gamma$ form the three biological contexts. A source-frozen predictor is fit at the source context and evaluated without re-training in the target context. Fig.~1B retains all 243 background rank trajectories and highlights a deterministic set of source-rank quantile exemplars for legibility; the full shortlist accounting retains all 243 labels.

The independent HepG2--Jurkat replication uses the Nadig data set \citep{nadig2025trade}. Its target-sum/$\log(1+x)$ preprocessing is fixed before evaluation. The primary analysis requires 20 cells per condition, uses five label-disjoint folds, three frozen ensemble members, and 30 resampling seeds. These are evaluation settings, not fitted outcomes.

\begin{table}[H]
\centering
\small
\caption{Eligibility conditions for the independent replication. The full registry, including rejected and incomplete candidates, remains a repository artifact rather than a manuscript table.}
\label{tab:registry}
\begin{tabular}{lll}
\toprule
Criterion & Requirement & Nadig evidence \\
\midrule
Shared perturbations & matched labels across contexts & 2,392 \\
Cell support & minimum 20 cells per condition & 87.2\% eligible \\
Matched controls & controls available in both contexts & 4,976 / 12,013 cells \\
Technical comparability & same library, RNA readout, and harvest design & supported \\
Major context confounding & no obvious context-linked batch confound & none identified \\
\bottomrule
\end{tabular}
\end{table}

\section{Measurement-adjusted reliability transport}
\label{app:estimators}

Pairwise risk comparisons retain ties as an explicit ordering state rather than coercing them into a strict direction. Separately, exact/tolerance ties are excluded from the strict-sign Beta posterior used for the stable-direction diagnostic. A strict direction is declared stable when the 95\% equal-tail Beta(1,1) posterior credible interval excludes $0.5$, with at least eight non-tie observations supporting the strict comparison; otherwise the pair remains unresolved.

We report the finite-sampling floor, the measurement-adjusted discrepancy, and a conservative joint-floor stress test separately. Effect-estimate intervals resample matched perturbation labels together with the declared measurement-seed and model-member resampling units. The decision-link association is descriptive over 18 directed source--target surface rows; its dependence-aware sensitivity resamples three unordered context-pair blocks while preserving both directions and all metric rows. No target labels or target-derived features enter the source-frozen predictor.

For a finite set of $M$ resampling replicates, let $Z_r$ denote the
tie-aware state for one perturbation pair in replicate $r$. The within-context
coincidence term is estimated with the unbiased order-2 U-statistic
\[
 \widehat W_U=\frac{1}{M(M-1)}\sum_{r\ne r'}\mathbf{1}\{Z_r=Z_{r'}\},
 \qquad \widehat D_{\rm within}=1-\widehat W_U.
\]
This correction removes the positive finite-replicate bias of the plug-in
coincidence estimate under a null context pair. All exact comparisons use a
fixed tolerance of $10^{-8}$. The strict-direction diagnostic uses a
Beta$(1,1)$ prior after excluding exact/tolerance ties; its posterior support
count is the number of remaining strict observations. The metric-dependence
diagnostic uses 15 prespecified discovery seeds and 15 held-out validation
seeds, while the primary effect estimates retain the declared 30-seed
resampling contract.

The component analysis uses the empirical pair-order law rather than the
posterior tail probability used for stable-state classification. For each
directed transfer and metric, both measurement replicates are retained and
independent 15-seed halves form a cross-product estimator. The finite-sample
report contains a tie component and a directional-bias component, compared
with the canonical per-seed U-corrected $D_{\rm adj}$ on the same eligible
universe. The population strength/flip identity is retained algebraically,
but its nonlinear empirical split is not treated as unbiased. At the
population level $0\leq D_{\rm adj}\leq 1$, but the unbiased
  finite-sample estimate $\widehat D_{\rm adj}$ is signed and is not clipped;
negative values indicate that the observed pair-state discrepancy is smaller
than the finite-replicate correction. In the deterministic tie-free limit,
each pair law is a point mass and $D_{\rm adj}$ reduces exactly to the
normalized Kendall inversion fraction $(1-\tau)/2$. This is a limit check,
not a substitute for the finite-measurement estimand.

\section{Extended transport evidence and independent replication}
\label{app:extended}

The full-surface transport diagnostics contain the shared labels available to each matched context pair and are reported in the figure-source tables. Aggregate $D_{\rm adj}$ estimates are formed from globally merged eligible per-transfer sets; shard-local pairwise summaries are never averaged. Same-context controls are near zero (0.001--0.005), showing that the signal is not a generic artifact of the pairwise estimator.

Table~\ref{tab:nadig} gives the matched-universe sensitivity of the independent Nadig replication at the two prespecified eligibility thresholds (minimum 20 versus 40 cells per condition). The six rows cover all three metrics at each threshold; the complete two-direction full-size surface remains in the figure source table. The values are presented as measurement-adjusted components rather than as conventional prediction accuracy, keeping the replication aligned with the main claim.

For universe accounting, the Nadig registry contains 2,392 shared labels;
2,086 labels per direction meet the primary minimum-20-cell eligibility rule,
and the matched sensitivity universes contain 1,255 labels at minimum 20 and
345 labels at minimum 40 after intersecting the two directions and the fixed
metric panel. These counts are registry properties, not model-selected
subsets.

\begin{table}[H]
\centering
\caption{Nadig matched-universe sensitivity: measurement-adjusted ordering disagreement ($D_{\rm adj}$; 90\% bootstrap interval) at two minimum-support thresholds.}
\label{tab:nadig}
\small
\begin{tabular}{lll}
\toprule
Universe & Metric & $D_{\rm adj}$ [90\% CI] \\
\midrule
min20 / 1255 & Delta cosine & 0.170 [0.152, 0.188] \\
min20 / 1255 & Systema centroid & 0.119 [0.092, 0.141] \\
min20 / 1255 & Absolute-effect rank & 0.108 [0.093, 0.124] \\
min40 / 345 & Delta cosine & 0.168 [0.128, 0.204] \\
min40 / 345 & Systema centroid & 0.085 [0.047, 0.124] \\
min40 / 345 & Absolute-effect rank & 0.073 [0.044, 0.103] \\
\bottomrule
\end{tabular}
\end{table}


\section{Measurement and decision robustness}
\label{app:robustness}

Depth analyses use matched pseudoreplicates at 10, 20, 40, 80, and 160 cells per perturbation--context condition, with the full-depth estimate as reference. Detectability is the first depth whose 90\% interval lower bound exceeds zero; resolution is the first depth within 10\% of the full-depth estimate.

The decision link is robust over 5\%, 10\%, 20\%, and 50\% shortlist budgets, using the random-reference regret scale of the main analysis and a bounded worst-oracle scale as a secondary diagnostic. Table~\ref{tab:decision-robustness} reports the descriptive surface association and unordered-pair/source-context leave-out ranges. These rows are a compact counterpart to the main-text decision figure, not a second set of decision claims.

\begin{table}[H]
\centering
\caption{Decision-link robustness across shortlist budgets and regret normalizations. The random-reference rows use the same canonical full-depth decision surface as Fig.~4B; the bounded worst-oracle rows are a secondary source-frozen diagnostic. The last column is the leave-one-unordered-context-pair-out range.}
\label{tab:decision-robustness}
\small
\begin{tabular}{llcc}
\toprule
Budget & Regret scale & $\rho$ & Unordered-pair LOPO range \\
\midrule
5\% & random-reference normalized regret & 0.769 & [0.748, 0.783] \\
5\% & worst-oracle-normalized regret & 0.701 & [0.657, 0.776] \\
10\% & random-reference normalized regret & 0.775 & [0.755, 0.811] \\
10\% & worst-oracle-normalized regret & 0.796 & [0.755, 0.818] \\
20\% & random-reference normalized regret & 0.886 & [0.860, 0.916] \\
20\% & worst-oracle-normalized regret & 0.806 & [0.741, 0.839] \\
50\% & random-reference normalized regret & 0.948 & [0.930, 0.958] \\
50\% & worst-oracle-normalized regret & 0.889 & [0.832, 0.902] \\
\bottomrule
\end{tabular}
\end{table}


\subsection{Transport decomposition and rank comparators}

The full-surface decomposition contains 18 directed transfer--metric rows on
the matched-fixed universe (139 or 186 perturbations, hence 9,591 or 17,205
unordered pairs). Table~\ref{tab:transport-decomposition} reports a
finite-measurement cross-product for the tie and directional-bias components,
and compares their sum with canonical per-seed U-corrected $D_{\rm adj}$.
The population strength/flip identity is retained algebraically, but its
nonlinear empirical split is not presented as unbiased. Table~\ref{tab:rank-comparator}
compares alternative rank summaries on the identical 18-row surface; these are
descriptive comparator analyses and do not change the primary estimand. The
depth-convergence ledger records each comparator, its same-context floor, and
$|T(n)-T(\mathrm{full})|$ for every matched depth.

\begin{table}[H]
\centering
\scriptsize
\caption{Finite-measurement cross-product decomposition on the matched-fixed full-depth surface. The empirical report contains the tie and directional-bias components; the population strength/flip identity is not treated as an unbiased finite-sample split.}
\label{tab:transport-decomposition}
\begin{tabular}{llrrrrr}
\toprule
Transfer & Metric & $D_{\rm tie}^{\rm CF}$ & $D_{\rm dir}^{\rm CF}$ & Sum & $D_{\rm adj}$ & Difference \\
\midrule
Co $\to$ Ctrl & Abs.-effect rank & 0.000 & 0.090 & 0.090 & 0.072 & 0.017 \\
Co $\to$ Ctrl & $\Delta$ cosine & 0.000 & 0.160 & 0.160 & 0.147 & 0.013 \\
Co $\to$ Ctrl & Systema & 0.008 & 0.083 & 0.091 & 0.073 & 0.018 \\
Co $\to$ IFN$\gamma$ & Abs.-effect rank & 0.000 & 0.083 & 0.083 & 0.067 & 0.016 \\
Co $\to$ IFN$\gamma$ & $\Delta$ cosine & 0.000 & 0.152 & 0.152 & 0.135 & 0.017 \\
Co $\to$ IFN$\gamma$ & Systema & 0.008 & 0.053 & 0.061 & 0.051 & 0.010 \\
Ctrl $\to$ Co & Abs.-effect rank & 0.000 & 0.101 & 0.101 & 0.085 & 0.017 \\
Ctrl $\to$ Co & $\Delta$ cosine & 0.000 & 0.181 & 0.181 & 0.164 & 0.018 \\
Ctrl $\to$ Co & Systema & 0.008 & 0.080 & 0.088 & 0.070 & 0.018 \\
Ctrl $\to$ IFN$\gamma$ & Abs.-effect rank & 0.000 & 0.109 & 0.109 & 0.090 & 0.019 \\
Ctrl $\to$ IFN$\gamma$ & $\Delta$ cosine & 0.000 & 0.156 & 0.156 & 0.134 & 0.022 \\
Ctrl $\to$ IFN$\gamma$ & Systema & 0.008 & 0.067 & 0.075 & 0.060 & 0.015 \\
IFN$\gamma$ $\to$ Co & Abs.-effect rank & 0.000 & 0.085 & 0.085 & 0.071 & 0.014 \\
IFN$\gamma$ $\to$ Co & $\Delta$ cosine & 0.000 & 0.181 & 0.181 & 0.164 & 0.016 \\
IFN$\gamma$ $\to$ Co & Systema & 0.008 & 0.050 & 0.058 & 0.044 & 0.014 \\
IFN$\gamma$ $\to$ Ctrl & Abs.-effect rank & 0.000 & 0.092 & 0.092 & 0.079 & 0.014 \\
IFN$\gamma$ $\to$ Ctrl & $\Delta$ cosine & 0.000 & 0.146 & 0.146 & 0.127 & 0.019 \\
IFN$\gamma$ $\to$ Ctrl & Systema & 0.008 & 0.072 & 0.080 & 0.063 & 0.017 \\
\bottomrule
\end{tabular}
\end{table}

\begin{table}[H]
\centering
\scriptsize
\caption{Rank-comparator benchmark on the same 18 directed transfer--metric rows. The final column is the descriptive Spearman association with normalized regret.}
\label{tab:rank-comparator}
\begin{tabular}{lrl}
\toprule
Comparator & $n$ & Spearman with regret \\
\midrule
Kendall $(1-\tau_b)/2$ & 18 & 0.87 \\
Spearman $(1-\rho)/2$ & 18 & 0.84 \\
Top-10\% Jaccard distance & 18 & 0.84 \\
Stable-pair inversion & 18 & 0.85 \\
$D_{\rm adj}$ & 18 & 0.78 \\
\bottomrule
\end{tabular}
\end{table}


\subsection{Finite-depth discrimination against a matched null}

The canonical summary field \texttt{measurement\_floor} is the within-context
term subtracted by the primary estimator; it is not itself a same-context
$D_{\rm adj}$ shift. We therefore recomputed a matched pseudo-context null at
each depth by splitting the 30 seeds into independent 15-seed halves, keeping
both measurement replicates in each half. The cross-context estimate uses both
held-out orientations, and the same-context null uses the two within-context
half comparisons. This produces 30 macro rows and 540 surface rows across the
five comparator families and six depths. Because the finite $D_{\rm adj}$ unit
is consequential, we also compared this flattened 15-seed$\times$2-replicate
construction with a seed-block-preserving audit that U-corrects each aligned
two-replicate seed block before averaging. Across the 108 $D_{\rm adj}$ null
rows, the current construction was negative in 108/108 rows (mean $-0.00673$),
whereas the seed-block audit was negative in 63/108 (mean $-0.000334$); the
maximum absolute difference was $0.0302$ and rank correlation was $0.074$.
The constructions therefore do not estimate the same finite-sample quantity.
The main-text Fig.~3C is explicitly labeled a split-half pair-law
diagnostic, while the canonical primary $D_{\rm adj}$ remains the per-seed
U-corrected summary. No negative estimate was clipped or selected away.
All comparator margins were positive
at every depth, but their magnitudes differ substantially; in particular,
top-10 Jaccard has the largest full-depth margin whereas $D_{\rm adj}$ is the
finite-measurement pair-state formulation, not a universally dominant rank
summary. The stable-inversion comparator has lower evaluable coverage at low
depth and is interpreted accordingly. In addition, we materialized a separate
matched-unit audit for all five comparator families. Each aligned
two-replicate seed block is collapsed to one block-level vector for the rank
comparators; $D_{\rm adj}$ is U-corrected within each aligned block and
averaged, and the stable-direction fraction uses the same 15 block-level
vectors with the frozen support rule. This audit has one row per surface,
depth, and comparator (540 rows) and is not substituted for the primary
flattened-partition estimate. All five matched-unit margins remain positive
in 108/108 rows; for $D_{\rm adj}$, 3/108 matched cross estimates and 63/108
matched same-context null estimates are negative, while the margin remains
positive in every row. We therefore retain the signed estimates and do not
interpret the finite-sample null sign as a biological effect.

\begin{table}[H]
\centering
\caption{Finite-measurement discrimination at full depth. Each comparator is
evaluated on the same 18 directed transfer--metric surfaces. The margin is the
cross-context estimate minus the matched same-context pseudo-context null; the
interval is a 90\% unordered-context-pair bootstrap interval. These results do
not rank $D_{\rm adj}$ as the universally best comparator.}
\label{tab:finite-measurement}
\small
\begin{tabular}{lrrrr}
\toprule
Comparator & Cross & Same-context null & Margin [90\% CI] & Recovery \\
\midrule
Kendall distance & 0.256 & 0.028 & 0.229 [0.216, 0.242] & 1.00 \\
Spearman distance & 0.177 & 0.003 & 0.174 [0.157, 0.191] & 1.00 \\
Top-10 Jaccard distance & 0.738 & 0.131 & 0.607 [0.594, 0.620] & 1.00 \\
Stable inversion & 0.147 & 0.000 & 0.147 [0.134, 0.160] & 1.00 \\
$D_{\rm adj}$ & 0.106 & $-0.005$ & 0.111 [0.105, 0.116] & 1.00 \\
\bottomrule
\end{tabular}
\end{table}


\subsection{Incremental decision value and empirical target-audit replay}

To distinguish the pooled decision association from a mean-fidelity scale
effect, we fit five prespecified nested models under leave-one-unordered-
context-pair-out
validation. The prespecified model family is evaluated under this outer split;
the nested contrasts are M1--M0, M2--M0 and M3--M1, while M4 is retained as an
exploratory decomposition-feature model because it is not a superset of M3.
All continuous features are standardized within the training fold.
The mean-shift model gives a modest OOF MAE improvement of
$\PTLIncrementalMeanDeltaMAE$ with an unordered-pair block-bootstrap 90\% interval
$\PTLIncrementalMeanDeltaMAECI$. The $D_{\rm adj}$ contrast beyond metric plus
mean shift is $\PTLIncrementalDAdjBeyondMeanDeltaMAE$
($\PTLIncrementalDAdjBeyondMeanDeltaMAECI$). These are uncertainty-aware
diagnostics, not causal effects. Table~\ref{tab:incremental-decision-value}
gives the complete OOF comparison and 2,000-draw unordered-pair block
bootstrap; no headline permutation $p$-value is reported for three blocks.

\begin{table}[H]
\centering
\tiny
\caption{Prespecified model comparison under leave-one-unordered-context-pair-out validation. $\Delta\mathrm{MAE}_{M0}=\mathrm{MAE}(M0)-\mathrm{MAE}(M)$ and $\Delta\mathrm{MAE}_{M1}=\mathrm{MAE}(M1)-\mathrm{MAE}(M)$, so positive values indicate improvement over the named reference. Intervals are 2,000 unordered-pair block-bootstrap diagnostics. M4 is exploratory because its decomposition features are not a superset of M3.}
\label{tab:incremental-decision-value}
\resizebox{\linewidth}{!}{%
\begin{tabular}{lrrrrrrr}
\toprule
Model & MAE & RMSE & OOF $\rho$ & $\Delta\mathrm{MAE}_{M0}$ & $\Delta\mathrm{MAE}_{M1}$ & 90\% interval & Validation \\
\midrule
M0 metric & 0.100 & 0.130 & 0.60 & 0.000 & -- & -- & ref. \\
M1 metric plus mean & 0.087 & 0.119 & 0.66 & 0.013 & 0.000 & [-0.001, 0.027] & LOPO (3) \\
M2 metric plus d adj & 0.105 & 0.132 & 0.64 & -0.005 & -- & [-0.019, 0.008] & LOPO (3) \\
M3 metric plus mean plus d adj & 0.092 & 0.121 & 0.62 & 0.008 & -0.005 & [-0.011, 0.026] & LOPO (3) \\
M4 metric plus mean plus decomposition (exploratory) & 0.106 & 0.137 & 0.66 & -0.007 & -0.019 & [-0.040, 0.027] & LOPO (3) \\
\bottomrule
\end{tabular}
}
\end{table}


The shortlist audit uses lower/upper target-risk intervals and the rules
\textsc{reuse} if $\mathrm{UB}_{\rm regret}\leq\epsilon$, \textsc{revise} if
$\mathrm{LB}_{\rm regret}>\epsilon$, and \textsc{unresolved} otherwise. The
offline replay uses $\epsilon=0.05$, gives every candidate a 10-cell target
pilot, and spends an additional target-cell budget through the 20/40/80/160
depth ladder. Policies see only the audit pool; the disjoint reference pool is
retained offline and the evaluation pool is used only for realized regret.
This is an empirical decision replay, not a certification framework or a
simultaneous confidence guarantee. Table~\ref{tab:target-audit} summarizes all
$\PTLTargetAuditRows$ rows; the row-level ledger is retained in the repository.

\begin{table}[H]
\centering
\scriptsize
\caption{Empirical target-audit replay. Every candidate receives a 10-cell pilot; the policy allocates additional target cells using only the audit pool. Values average the four decision budgets.}
\label{tab:target-audit}
\begin{tabular}{lrrrr}
\toprule
Policy & Mean upgraded & False REUSE & Correct REUSE & Correct REVISE \\
\midrule
pilot only & 0.0 & 0.000 & 0.333 & 0.028 \\
ptl boundary & 55.7 & 0.097 & 0.444 & 0.056 \\
random extra depth & 32.6 & 0.000 & 0.333 & 0.069 \\
source uncertainty & 33.0 & 0.000 & 0.333 & 0.125 \\
uniform depth & 144.9 & 0.000 & 0.333 & 0.056 \\
\bottomrule
\end{tabular}
\end{table}


\subsection{Regret-bound-critical active-audit redesign}

As a final prespecified attempt to improve the active audit, we allocated
additional cells to candidates active in the regret upper bound: the highest
upper-bound items inside the source shortlist and the lowest lower-bound items
outside it, followed by interval width. At matched actual additional-cell
cost, false reuse was eliminated relative to the legacy PTL boundary, but the
policy did not exceed the source-uncertainty baseline in correct revision. The
realized evaluation-pool regret is invariant under the frozen-shortlist replay
by construction and is reported only as a policy-invariant reference. The
kill-switch therefore demotes this policy to a negative supplementary result
rather than a main-text contribution.

\begin{table}[H]
\centering
\caption{One-last active-audit redesign at matched actual additional-cell
cost. The regret-bound-critical policy reduces false reuse relative to the
legacy PTL boundary, but does not exceed the source-uncertainty baseline in
correct revision at matched cost. Realized evaluation-pool regret is invariant
under this frozen-shortlist replay by construction and is reported only as a
policy-invariant reference. The prespecified decision rule therefore retains it
as a negative supplementary result.}
\label{tab:regret-bound-critical}
\small
\begin{tabular}{lrrrr}
\toprule
Policy & False reuse & Correct revise & True regret & Normalized regret \\
\midrule
Regret-bound critical & 0.000 & 0.083 & 0.040 & 0.378 \\
Random extra depth & 0.000 & 0.069 & 0.040 & 0.378 \\
Uniform depth & 0.000 & 0.056 & 0.040 & 0.378 \\
Source uncertainty & 0.000 & 0.125 & 0.040 & 0.378 \\
Legacy PTL boundary & 0.097 & 0.056 & 0.040 & 0.378 \\
\bottomrule
\end{tabular}
\end{table}


\section{Can reordering be anticipated from source-only information?}
\label{app:prospective}

This analysis asks a narrower prospective question than the transport study: can a source-only feature ladder anticipate future reordering burden? The legal feature sets are cumulative uncertainty ($U$), geometry ($U+G$), support ($U+G+S$), and novelty ($U+G+S+N$). Target outcomes, target risk, and the measurement-adjusted transport value are excluded by an information firewall. The source-only task-specific usable universe contains approximately 28--44 labels; the 243-label universe belongs to the full-size shortlist surface.

The evaluation contains 18 transfer-by-metric tasks for each of continuous burden and high-burden classification. It uses ten repeats of five label-disjoint folds with nested ridge selection. Every source row records \texttt{target\_feature\_leakage=False}. Supplementary Fig.~\ref{fig:source-only} shows held-out task-level performance directly; it intentionally does not recycle the earlier deployment-gate analysis or an unreconstructible chance baseline.

\begin{figure}[t]
\centering
\includegraphics[width=\textwidth]{../../results/figures/reliability_transportability/reliability_transportability_supp_fig1_source_only_forecasting.pdf}
\caption{Source-only prospective boundary. (A) The information firewall separates source-observable features from target outcomes. (B) Held-out Spearman correlation for continuous reordering burden. (C) Held-out AUPRC for high-burden classification. Each point is one transfer-by-metric task; thick marks summarize the median and interquartile range.}
\label{fig:source-only}
\end{figure}

\section{Additional robustness and diagnostic analyses}
\label{app:reviewer-proofing}

Supplementary Fig.~\ref{fig:metric-stratified} checks whether the primary
decision association is explained only by the different scales of the three
metrics. Panel A shows the fixed 18-point surface; Panel B ranks $D_{\rm adj}$
and regret within each metric before pooling. The pooled within-metric
association is descriptive ($\rho=\PTLWithinMetricRho$). With only three
unordered context-pair blocks, the percentile 90\% block-bootstrap interval
is a diagnostic, $\PTLWithinMetricRhoCI$ from 2,000 draws; leave-one-pair-out
estimates range from $0.00$ to $0.47$. The six-point metric-specific summaries
are shown without a headline significance claim.

\begin{figure}[H]
\centering
\includegraphics[width=\textwidth]{../../results/figures/reliability_transportability/reliability_transportability_supp_fig2_metric_stratified_decision.pdf}
\caption{Metric-stratified decision association. The 18 fixed full-depth,
10\% random-reference decision rows are shown in their native scale (A) and
after within-metric rank transformation (B). The pooled within-metric interval
is reported in the text and source table rather than overlaid on the rank
points. Colors identify metrics; no target-informed feature enters the
source-frozen predictor.}
\label{fig:metric-stratified}
\end{figure}

\FloatBarrier

Supplementary Fig.~\ref{fig:mean-fidelity} computes the absolute change in
mean metric risk on the exact matched perturbation universe used by
$D_{\rm adj}$. It is a metric-matched control for the statement that a change
in average fidelity and a change in reliability ordering are distinct
quantities; marker shape identifies directed transfer and color identifies the
metric. Across the 18 rows, the mean-risk shift has Spearman $\rho=\PTLMeanFidelityShiftRhoD$
with $D_{\rm adj}$ and $\rho=\PTLMeanFidelityShiftRhoRegret$ with normalized
regret; these are descriptive matched-universe associations, not a causal
decomposition.

\begin{figure}[t]
\centering
\includegraphics[width=\textwidth]{../../results/figures/reliability_transportability/reliability_transportability_supp_fig3_mean_fidelity_control.pdf}
\caption{Metric-matched mean-fidelity control. Each point is one of 18
directed transfer--metric rows. The horizontal coordinate is the absolute
source-to-target mean-risk shift computed on the exact fixed label universe;
the vertical coordinate is $D_{\rm adj}$ (A) or the corresponding normalized
decision regret at 10\% budget (B).}
\label{fig:mean-fidelity}
\end{figure}

\FloatBarrier

Supplementary Fig.~\ref{fig:estimand-validation} validates the finite-sample
estimand in synthetic data with a known null and a known alternative. The
stored trial-level distributions show the positive plug-in bias under the
null and the centering of the U-corrected estimator near the population
identity.

\begin{figure}[t]
\centering
\includegraphics[width=\textwidth]{../../results/figures/reliability_transportability/reliability_transportability_supp_fig4_estimand_validation.pdf}
\caption{Finite-sample validation of the ordering estimand. Trial-level
sampling distributions compare the plug-in/V estimator with the order-2
U-corrected estimator under a finite-sample null and a known alternative;
vertical references mark the population identity.}
\label{fig:estimand-validation}
\end{figure}

\section{External predictor context under matched information constraints}
\label{app:external}

To broaden the model-family context without changing the primary estimand, we
also record the completed formal-v2 external GEARS contract~\citep{roohani2023gears}
on three public
cohorts. These canonical runs use condition-disjoint holdout, zero perturbation
overlap, a dataset-native frozen 4,096-gene panel, three seeds and 20 training
epochs.
The two larger cohorts include held-out validation predictions; the K562 GWPS
cohort has only 64 test rows and is therefore retained as a test-only audit
surface. Supplementary Table~\ref{tab:external-gears} reports the complete
dataset-level summary. These results characterize conventional response
prediction under external cohorts; they neither establish predictor-family
invariance nor enter the Frangieh context-shift claim lock.

\begin{table}[H]
\centering
\small
\setlength{\tabcolsep}{4pt}
\caption{\textbf{Scope-limited external GEARS audit on public cohorts.}
Each row aggregates three CUDA runs (20 epochs; seeds 0--2) under a
random-condition-holdout split with zero perturbation overlap and a
dataset-native frozen 4,096-gene panel. Cosine similarity is reported on
non-control test cells (higher is better); MSE and MAE use the same cells
(lower is better). ``Validation'' counts runs with held-out validation
predictions. These conventional response-prediction results are not pooled
with the Frangieh source-to-target reliability-transport estimand.}
\label{tab:external-gears}
\begin{tabular}{@{}lrrrrll@{}}
\toprule
Dataset & Runs & Test cells & Cosine & MSE & MAE & Validation \\
\midrule
Norman filtered & 3 & 1,069 & 0.159 $\pm$ 0.058 & 0.01258 & 0.07552 & 3/3 \\
Replogle RPE1 & 3 & 983 & 0.136 $\pm$ 0.018 & 0.02245 & 0.11218 & 3/3 \\
Replogle K562 GWPS & 3 & 64 & 0.069 $\pm$ 0.021 & 0.02716 & 0.12828 & 0/3$^{\dagger}$ \\
\bottomrule
\end{tabular}
\par\noindent
\raggedright\footnotesize $^{\dagger}$Test-only external audit surface; no validation vector was available.
\end{table}


An earlier one-epoch GEARS pilot used the same dataset/seed labels and frozen
panels but produced a different set of predictions. It is retained only in the
provenance audit (\texttt{gears\_duplicate\_run\_audit.csv}), is not pooled as
an additional replicate, and is excluded from both external tables; the
formal-v2 family above is the sole canonical GEARS source.

Using the same frozen panels and condition-disjoint split construction, we also completed five
non-GEARS controls (45 runs total): control mean, global perturbation delta,
perturbation mean delta, context-nearest-neighbour delta, and ridge
regression. Together with GEARS this gives 54 comparable runs across six
audited model/control families. Supplementary Table~\ref{tab:external-model-summary} reports
equal-run aggregate metrics; the complete run-level CSV is retained so that
the small-cohort sensitivity and model-specific variation remain auditable.
Because the holdout is perturbation-disjoint, the context-$k$NN control falls
back to the global mean on unseen test perturbations; we retain this row as an
explicit protocol control rather than presenting it as successful neighbour
transfer. These controls broaden the external benchmark context without
changing the source-frozen Frangieh estimand or implying a legal second
predictor family on that primary surface.

\begin{table}[H]
\centering
\small
\setlength{\tabcolsep}{3.5pt}
\caption{\textbf{Comparable external model-family audit on public cohorts.}
Each entry is an equal-weight mean across nine runs (three cohorts $\times$
three seeds) on the corresponding cohort-specific frozen 4,096-gene panel and
condition-disjoint
holdout construction. Per-seed test-cell counts match across the audited
model outputs. Cosine is higher-is-better; MSE and MAE are lower-is-better.
The GEARS rows are the verified CUDA runs in Table~\ref{tab:external-gears};
the remaining rows are deterministic or fitted baselines evaluated with the
same split construction. These summaries are conventional response-prediction
controls and are not pooled with the Frangieh reliability-transport claim.}
\label{tab:external-model-summary}
\begin{tabular}{@{}lrrrrr@{}}
\toprule
Model family & Runs & Cohorts & Cosine & MSE & MAE \\
\midrule
Ridge regression & 9 & 3 & 0.233 & 0.01887 & 0.09970 \\
Perturbation mean $\Delta$ & 9 & 3 & 0.202 & 0.01705 & 0.09548 \\
Global mean $\Delta$ & 9 & 3 & 0.185 & 0.01737 & 0.09624 \\
Context $k$NN $\Delta$ (fallback) & 9 & 3 & 0.185 & 0.01737 & 0.09624 \\
GEARS & 9 & 3 & 0.122 & 0.02073 & 0.10533 \\
Control mean & 9 & 3 & 0.000 & 0.01823 & 0.09852 \\
\bottomrule
\end{tabular}
\par\noindent
\raggedright\footnotesize Equal-run means are not test-cell-weighted; the K562
GWPS cohort contributes only 64 non-control test cells per three-run cohort.
Under the perturbation-disjoint holdout, the $k$NN control falls back to the
global mean for unseen test perturbations; it is retained as a protocol
control, not as evidence of context-neighbour transfer. Full run-level values
and output paths are in
\texttt{results/tables/external\_condition\_holdout\_shared\_panel\_metrics.csv}.
\end{table}


The broader predictor-family plan was also audited against the same
source-only contract. Supplementary Table~\ref{tab:predictor-compatibility}
records the executable or semantic status of the additional families, including
scGPT~\citep{cui2024scgpt}, TxPert~\citep{wenkel2026txpert}, CPA~\citep{lotfollahi2023cpa},
scGen~\citep{lotfollahi2019scgen} and CellOT~\citep{bunne2023cellot}. These
rows are deliberately not treated as performance comparisons: a package,
checkpoint, or context-mismatched inference run is not evidence for the
Frangieh transport claim. The corresponding machine-readable ledgers are
\texttt{predictor\_family\_protocol.csv} and
\texttt{modern\_predictor\_feasibility.json}.

\begin{table}[H]
\centering
\small
\setlength{\tabcolsep}{3.5pt}
\caption{\textbf{Predictor-family compatibility audit.} The audit records
whether a public implementation can satisfy the manuscript's frozen
source-only prediction contract: training uses source information only,
target outcomes are unseen until evaluation, and the output is a held-out
vector on the declared gene panel. Only the canonical formal-v2 GEARS cohort
runs enter Supplementary Tables~\ref{tab:external-gears}--\ref{tab:external-model-summary}; the other rows are feasibility or semantic-contract
records, not performance results.}
\label{tab:predictor-compatibility}
\begin{tabular}{@{}p{0.19\linewidth}p{0.24\linewidth}p{0.49\linewidth}@{}}
\toprule
Family & Current evidence & Manuscript decision \\
\midrule
GEARS & Formal-v2 CUDA runs completed on three public cohorts; 20 epochs; native uncertainty retained & Included only as the scope-limited external cohort audit in Supplementary Tables~\ref{tab:external-gears}--\ref{tab:external-model-summary}; not pooled with the Frangieh estimand. \\
State & Official checkout and preprocessing verified; its set-transition output does not provide the required frozen global five-fold source-only vector and target-risk bridge & Outside the declared comparison; no score reported. \\
TxPert & Official cached inference is reproducible for K562, but the available checkpoint/data do not provide the frozen Nadig HepG2/Jurkat source-only vector & Blocked by context mismatch; K562 performance is not relabelled as cross-cell transfer. \\
scGPT & A source-only held-out vector on the declared Frangieh panel is not available under the prespecified contract & No score is claimed; the family remains outside the transport comparison. \\
biolord & Source-only adapter and pinned compatible checkpoint are not available in the current workspace & Compatibility pending; no score is claimed. \\
PerturbNet & Source-only adapter and pinned compatible checkpoint are not available in the current workspace & Compatibility pending; no score is claimed. \\
Scouter & Source-only adapter and pinned compatible checkpoint are not available in the current workspace & Compatibility pending; no score is claimed. \\
CPA-ext / scGen / CellOT & Rescue-layer implementations and weights are not information-matched to the declared source-only vector contract & Deferred to a separate context-aware rescue audit; no score is claimed. \\
\bottomrule
\end{tabular}
\end{table}


\subsection{Source-only predictor-family sensitivity on the Frangieh surface}

To test whether the transport profile is tied to the primary bilinear-ridge
parameterization, we materialized a source-only RBF kernel-ridge family on the
same Frangieh labels, gene panel, three metrics, six directed transfers and
full-size minimum-20 measurement surface. The target outcomes were loaded
only after source-only prediction materialization; they were not visible during
fit, penalty tuning or family selection. The initial penalty grid was extended
because all 15 source-by-outer-fold choices selected its largest value. The
resulting RBF sensitivity preserves the positive adjusted-discrepancy interval
on all 18 matched rows and has descriptive $D_{\rm adj}$ profile Spearman
$\rho=0.983$ with the bilinear-ridge family; the within-metric profile range is
$0.829$--$1.000$ and the mean row-wise prediction-vector cosine is $0.733$.
Point and interval sign agreement holds on all 18 rows. This is a matched-surface sensitivity, not evidence
of predictor-family equivalence, superiority, causality or broader
generalization.

\begin{table}[H]
\centering
\scriptsize
\setlength{\tabcolsep}{2pt}
\caption{\textbf{Source-only predictor-family sensitivity on the matched
Frangieh surface.} The RBF kernel-ridge family was tuned using source-only
inner cross-validation; because all 15 initial source-by-outer-fold choices
hit the initial largest penalty, the grid was extended to
$\{30,100,300,1000\}$. The comparison uses the same full-size minimum-20
label universes, frozen metrics, transfer directions and measurement-floor
estimator. Source OOF fidelity is shown as the macro mean for delta cosine,
Systema centroid accuracy and absolute-effect rank agreement, respectively.
Concordance is descriptive and does not establish predictor-family
  equivalence or superiority. For the RBF family, the matched profile
  concordance is $\rho=0.983$, with within-metric range $0.829$--$1.000$
  and mean row-wise prediction-vector cosine $0.733$.}
\label{tab:predictor-family-sensitivity}
\resizebox{\linewidth}{!}{%
\begin{tabular}{@{}p{0.16\linewidth}p{0.22\linewidth}p{0.09\linewidth}p{0.22\linewidth}p{0.22\linewidth}@{}}
\toprule
Family & Source-only tuning & Surfaces & Source OOF fidelity ($\Delta$ / Sys / rank) & Matched transport concordance \\
\midrule
Bilinear ridge & Registered source-inner CV & 18 & 0.240 / 0.499 / 0.299 & Reference family \\
RBF kernel ridge & 15/15 hit $\alpha=10$; extended to $\alpha=30$ & 18 & 0.257 / 0.498 / 0.310 & $D_{\rm adj}$ profile $\rho=0.983$; point and CI sign agreement 18/18 \\
\bottomrule
\end{tabular}
}
\end{table}


\subsection{Provenance limits on a four-context expansion}

We audited the exact provenance and support before adding a broader
HepG2--Jurkat--K562--RPE1 transport map. HepG2 and Jurkat are both day-7
harvests, but the available Replogle K562 and RPE1 exports are day 6 and day 7,
respectively, with no timepoint field in the local cell-level exports. The
harvest difference is therefore inseparable from the intended context
contrast. Although the four-way label intersection and minimum-support counts
are large, the current 1,536-gene panel is missing 185 genes in at least one
context. We consequently do not present a clean four-context biological map;
the audit is a provenance boundary, not a negative biological result. The
underlying public-cohort provenance follows the Nadig atlas and the Replogle
genome-scale Perturb-seq resource \citep{nadig2025trade,replogle2022genome}.

\begin{table}[H]
\centering
\caption{Provenance audit for a possible four-context expansion. The
HepG2--Jurkat pair is a supported independent context comparison, but the
available K562 and RPE1 exports differ in harvest day and cannot support a
clean four-context biological map under the current contract.}
\label{tab:context-audit}
\small
\begin{tabular}{lll}
\toprule
Criterion & Result & Consequence \\
\midrule
HepG2--Jurkat harvest & Both day 7 & Supported independent pair \\
K562--RPE1 harvest & Day 6 vs day 7 & Context--time confounded \\
Four-way shared labels & 2,053 & Label support is adequate \\
Minimum-20 support & 1,688 labels & Support is adequate for profiling \\
Current 1,536-gene panel & 1,351 present in all four & Panel cannot be reused verbatim \\
Four-context map & Blocked & No clean biological claim \\
\bottomrule
\end{tabular}
\end{table}


\FloatBarrier

\section{Reproducibility and scientific regeneration}
\label{app:reproducibility}

The scientific regeneration path starts from the canonical manifest tables, fixed splits and resampling settings, and rebuilds the four main figures, Supplementary Figs.~S1--S4, and the compact supplementary tables. The anonymous repository documents the executable commands and data lineage. Large raw H5AD/NPZ inputs remain external to the manuscript tree; the reviewable CSV/JSON products expose the values used in figures and tables.

\paragraph{Code and data availability.}
The anonymized core release is available at
\href{https://github.com/Neabigmo/PTL/tree/iclr2027-anonymous-core}{the
anonymous core branch} and contains
the source code, repository-relative configurations, manuscript sources,
publication figures, and compact reviewable data products. Raw single-cell
matrices are not redistributed. The Frangieh RNA screen and the Nadig
HepG2--Jurkat screens are available from the public scPerturb v1.4 Zenodo
record (\href{https://zenodo.org/records/13350497}{Zenodo}); the repository README
documents the relevant file names and regeneration commands.

The scope-limited external GEARS audit and its matched baseline controls are
summarized in Supplementary Tables~\ref{tab:external-gears} and
\ref{tab:external-model-summary}; the broader predictor-family compatibility
ledger is summarized in Table~\ref{tab:predictor-compatibility}. Broader
condition-holdout analyses, old failure-anatomy graphics, and the earlier
broad source-only deployment gate remain repository-only. Runtime compatibility
notes for predictor families that do not satisfy the frozen source-only
contract are likewise retained as an engineering ledger, not presented as
primary scientific evidence.

\bibliographystyle{../../third_party/iclr2027/iclr2027_conference}
\bibliography{../../manuscript/references}
\end{document}
